Across language models from the Gemma, Qwen, and OLMo families, we find a model-specific
linear residual-stream direction that separates human from AI self-representation.
We identify this as the human-AI axis, a direction that can
be causally manipulated by steering the model's position along it. A human-versus-AI
separation is already present in an untuned Gemma 3 base model, and across OLMo
2 checkpoints we find that the axis forms within the first 3.7\% of pretraining,
long before instruction tuning. Steering along the axis reliably flips a
model's stated identity, even on questions held out from axis extraction, and
we find the same effect in three language models across the Gemma and Qwen model
families. Unlike prompting, steering intervenes directly on this internal
representation while holding the prompt fixed, letting us test
its causal role rather than just its correlation with what the model says.
The representation is also sensitive to context, as a user claiming to
be an AI shifts the model's own self-representation significantly toward
being an AI on all three models, with no steering. Although steering along the
axis shifts a model's own self-report far more than its behavior, steering
toward a human self-representation produces safety-adverse behavioral changes. For example, it raises
the model's willingness to cause harm significantly on Gemma 3 4B and Qwen 2.5
7B, and directionally on Gemma 3 12B. On Gemma 3 4B, human-ward steering causes
prompt-injection execution to rise 16-fold while 94\% of rollouts stay
on-task. Writing register such as warm versus formal phrasing shifts the
raw projections more than identity does, but steering with a register-matched direction
does not reproduce the behavioral effects on Gemma 3 4B, leaving identity as the
best-supported interpretation of what steering changes. The human-AI axis is therefore
a safety-relevant inference-time variable whose behavioral effects are not captured
by self-report alone.
\iffinalcopy Code and data are available at \url{https://github.com/asherk7/human-ai}.\fi
